% !TeX root = RJwrapper.tex
\title{Blocking: An R Package for Blocking of Records for Record Linkage and Deduplication}


\author{by Maciej Beręsewicz and Adam Struzik}

\maketitle

\abstract{%
An abstract of less than 250 words.
}

\section{Introduction}\label{introduction}

Interactive data graphics provides plots that allow users to interact them. One of the most basic types of interaction is through tooltips, where users are provided additional information about elements in the plot by moving the cursor over the plot.

This paper will first review some R packages on interactive graphics and their tooltip implementations. A new package \CRANpkg{ToOoOlTiPs} that provides customized tooltips for plot, is introduced. Some example plots will then be given to showcase how these tooltips help users to better read the graphics.

\section{Background}\label{background}

Some packages on interactive graphics include \CRANpkg{plotly} \citep{plotly} that interfaces with Javascript for web-based interactive graphics, \CRANpkg{crosstalk} \citep{crosstalk} that specializes cross-linking elements across individual graphics. The recent R Journal paper \CRANpkg{tsibbletalk} \citep{RJ-2021-050} provides a good example of including interactive graphics into an article for the journal. It has both a set of linked plots, and also an animated gif example, illustrating linking between time series plots and feature summaries.

\section{\texorpdfstring{Blocking of records using \texttt{blocking} function}{Blocking of records using blocking function}}\label{blocking-of-records-using-blocking-function}

\section{Integration with existing packages}\label{integration-with-existing-packages}

\section{Case study}\label{case-study}

\subsection{Record linkage example}\label{record-linkage-example}

Let us first load the required packages.

\begin{verbatim}
library(blocking)
library(data.table)
\end{verbatim}

We will demonstrate the use of \texttt{blocking} function for record linkage with the \texttt{foreigners} dataset included in the package. This fictional representation of the foreign population in Poland was generated based on publicly available information, preserving the distributions from administrative registers. It contains 110,000 rows with 100,000 entities. Each row represents one record, with the following columns:

\begin{itemize}
\tightlist
\item
  \texttt{fname} -- first name,
\item
  \texttt{sname} -- second name,
\item
  \texttt{surname} -- surname,
\item
  \texttt{date} -- date of birth,
\item
  \texttt{region} -- region (county),
\item
  \texttt{country} -- country,
\item
  \texttt{true\_id} -- person ID.
\end{itemize}

\begin{verbatim}
data(foreigners)
head(foreigners)
\end{verbatim}

\begin{verbatim}
#>     fname  sname    surname       date region country true_id
#>    <char> <char>     <char>     <char> <char>  <char>   <num>
#> 1:   emin            imanov 1998/02/05            031       0
#> 2: nurlan        suleymanli 2000/08/01            031       1
#> 3:   amio        maharrsmov 1939/03/08            031       2
#> 4:   amik        maharramof 1939/03/08            031       2
#> 5:   amil        maharramov 1993/03/08            031       2
#> 6:  gadir        jahangirov 1991/08/29            031       3
\end{verbatim}

We split the dataset into two separate files: one containing the first appearance of each entity in the \texttt{foreigners} dataset, and the other containing its subsequent appearances.

\begin{verbatim}
foreigners_1 <- foreigners[!duplicated(foreigners$true_id), ]
foreigners_2 <- foreigners[duplicated(foreigners$true_id), ]
\end{verbatim}

Now in both datasets we remove slashes from the \texttt{date} column and create a new string column that concatenates the information from all columns (excluding \texttt{true\_id}) in each row.

\begin{verbatim}
foreigners_1[, date := gsub("/", "", date)]
foreigners_1[, txt := paste0(fname, sname, surname, date, region, country)]
foreigners_2[, date := gsub("/", "", date)]
foreigners_2[, txt := paste0(fname, sname, surname, date, region, country)]
head(foreigners_1)
\end{verbatim}

\begin{verbatim}
#>     fname  sname    surname     date region country true_id
#>    <char> <char>     <char>   <char> <char>  <char>   <num>
#> 1:   emin            imanov 19980205            031       0
#> 2: nurlan        suleymanli 20000801            031       1
#> 3:   amio        maharrsmov 19390308            031       2
#> 4:  gadir        jahangirov 19910829            031       3
#> 5:   zaur         bayramova 19961006  01261     031       4
#> 6:   asif          mammadov 19970726            031       5
#>                              txt
#>                           <char>
#> 1:         eminimanov19980205031
#> 2:   nurlansuleymanli20000801031
#> 3:     amiomaharrsmov19390308031
#> 4:    gadirjahangirov19910829031
#> 5: zaurbayramova1996100601261031
#> 6:       asifmammadov19970726031
\end{verbatim}

We use the newly created columns in the \texttt{blocking} function, which relies on the default \CRANpkg{rnndescent} (Nearest Neighbor Descent) algorithm based on cosine distance. Additionally, we set \texttt{verbose\ =\ 1} to monitor progress.

\begin{verbatim}
set.seed(2025)
result_reclin <- blocking(x = foreigners_1$txt, 
                          y = foreigners_2$txt, 
                          verbose = 1, 
                          n_threads = 4)
\end{verbatim}

\begin{verbatim}
#> ===== creating tokens =====
#> ===== starting search (nnd, x, y: 100000, 10000, t: 1232) =====
#> ===== creating graph =====
\end{verbatim}

Now we examine the results of record linkage.

\begin{itemize}
\tightlist
\item
  We have created 6464 blocks.
\item
  The blocking process utilized 1232 columns (2 character shingles).
\item
  We have 3,916 blocks of 2 elements, 1,604 blocks of 3 elements,\ldots, 2 blocks of 7 elements.
\end{itemize}

\begin{verbatim}
result_reclin
\end{verbatim}

\begin{verbatim}
#> ========================================================
#> Blocking based on the nnd method.
#> Number of blocks: 6464.
#> Number of columns used for blocking: 1232.
#> Reduction ratio: 0.9999.
#> ========================================================
#> Distribution of the size of the blocks:
#>    2    3    4    5    6    7 
#> 3912 1597  931   21    1    2
\end{verbatim}

Structure of the object is as follows:

\begin{itemize}
\tightlist
\item
  \texttt{result} -- a \texttt{data.table} with identifiers and block IDs,
\item
  \texttt{method} -- name of the ANN algorithm used,
\item
  \texttt{deduplication} -- whether deduplication was applied,
\item
  \texttt{representation} -- whether shingles or vectors were used,
\item
  \texttt{metrics} -- metrics for quality assessment (here \texttt{NULL}),
\item
  \texttt{confusion} -- confusion matrix (here \texttt{NULL}),
\item
  \texttt{colnames} -- column names used for the comparison,
\item
  \texttt{graph} -- an \CRANpkg{igraph} object, mainly for visualization (here \texttt{NULL}).
\end{itemize}

\begin{verbatim}
str(result_reclin, 1)
\end{verbatim}

\begin{verbatim}
#> List of 8
#>  $ result        :Classes 'data.table' and 'data.frame': 10000 obs. of  4 variables:
#>   ..- attr(*, ".internal.selfref")=<externalptr> 
#>  $ method        : chr "nnd"
#>  $ deduplication : logi FALSE
#>  $ representation: chr "shingles"
#>  $ metrics       : NULL
#>  $ confusion     : NULL
#>  $ colnames      : chr [1:1232] "0a" "0b" "0c" "0m" ...
#>  $ graph         : NULL
#>  - attr(*, "class")= chr "blocking"
\end{verbatim}

The resulting \texttt{data.table} has four columns:

\begin{itemize}
\tightlist
\item
  \texttt{x} -- reference dataset (i.e.~\texttt{foreigners\_1}) -- this may not contain all units of \texttt{foreigners\_1},
\item
  \texttt{y} -- query (each row of \texttt{foreigners\_2}) -- this may not contain all units of \texttt{foreigners\_2},
\item
  \texttt{block} -- block ID,
\item
  \texttt{dist} -- distance between objects.
\end{itemize}

\begin{verbatim}
head(result_reclin$result)
\end{verbatim}

\begin{verbatim}
#>        x     y block      dist
#>    <int> <int> <num>     <num>
#> 1:     3     1     1 0.2216882
#> 2:     3     2     1 0.2122737
#> 3:    21     3     2 0.1172652
#> 4:    57     4     3 0.1863238
#> 5:    57     5     3 0.1379310
#> 6:    61     6     4 0.2307692
\end{verbatim}

Let's examine the first pair. Obviously, there are typos in the \texttt{fname} and \texttt{surname}. Nevertheless, the pair appears to be a match.

\begin{verbatim}
cbind(t(foreigners_1[3, 1:6]), t(foreigners_2[1, 1:6]))
\end{verbatim}

\begin{verbatim}
#>         [,1]         [,2]        
#> fname   "amio"       "amik"      
#> sname   ""           ""          
#> surname "maharrsmov" "maharramof"
#> date    "19390308"   "19390308"  
#> region  ""           ""          
#> country "031"        "031"
\end{verbatim}

Now we use the \texttt{true\_id} values to evaluate our approach.

\begin{verbatim}
matches <- merge(x = foreigners_1[, .(x = 1:.N, true_id)],
                 y = foreigners_2[, .(y = 1:.N, true_id)],
                 by = "true_id")
matches[, block := rleid(x)]
head(matches)
\end{verbatim}

\begin{verbatim}
#> Key: <true_id>
#>    true_id     x     y block
#>      <num> <int> <int> <int>
#> 1:       2     3     1     1
#> 2:       2     3     2     1
#> 3:      20    21     3     2
#> 4:      56    57     4     3
#> 5:      56    57     5     3
#> 6:      60    61     6     4
\end{verbatim}

We have 10,000 matched pairs. We use the \texttt{true\_blocks} parameter in the \texttt{blocking} function to specify the true block assignments. We obtain the quality metrics for the assessment of record linkage.

\begin{verbatim}
result_2_reclin <- blocking(x = foreigners_1$txt, 
                            y = foreigners_2$txt, 
                            verbose = 1, n_threads = 4, 
                            true_blocks = matches[, .(x, y, block)])
\end{verbatim}

\begin{verbatim}
#> ===== creating tokens =====
#> ===== starting search (nnd, x, y: 100000, 10000, t: 1232) =====
#> ===== creating graph =====
\end{verbatim}

\begin{verbatim}
result_2_reclin
\end{verbatim}

\begin{verbatim}
#> ========================================================
#> Blocking based on the nnd method.
#> Number of blocks: 6470.
#> Number of columns used for blocking: 1232.
#> Reduction ratio: 0.9999.
#> ========================================================
#> Distribution of the size of the blocks:
#>    2    3    4    5    6    7 
#> 3923 1592  932   20    1    2 
#> ========================================================
#> Evaluation metrics (standard):
#>      recall   precision         fpr         fnr    accuracy specificity 
#>     96.8738     78.7100      0.0038      3.1262     99.9957     99.9962 
#>    f1_score 
#>     86.8524
\end{verbatim}

\section{\texorpdfstring{Customizing tooltip design with \pkg{ToOoOlTiPs}}{Customizing tooltip design with }}\label{customizing-tooltip-design-with}

\pkg{ToOoOlTiPs} is a packages for customizing tooltips in interactive graphics, it features these possibilities.

\section{A gallery of tooltips examples}\label{a-gallery-of-tooltips-examples}

The \CRANpkg{palmerpenguins} data \citep{palmerpenguins} features three penguin species which has a lovely illustration by Alison Horst in Figure \ref{fig:penguins-alison}.

\begin{figure}
\includegraphics[width=1\linewidth,height=0.3\textheight,alt={A picture of three different penguins with their species: Chinstrap, Gentoo, and Adelie. }]{figures/penguins} \caption{Artwork by \@allison\_horst}\label{fig:penguins-alison}
\end{figure}

Table \ref{tab:penguins-tab-static} prints at the first few rows of the \texttt{penguins} data:

\begin{table}
\centering
\caption{\label{tab:penguins-tab-static}A basic table}
\centering
\fontsize{7}{9}\selectfont
\begin{tabular}[t]{l|l|r|r|r|r|l|r}
\hline
species & island & bill\_length\_mm & bill\_depth\_mm & flipper\_length\_mm & body\_mass\_g & sex & year\\
\hline
Adelie & Torgersen & 39.1 & 18.7 & 181 & 3750 & male & 2007\\
\hline
Adelie & Torgersen & 39.5 & 17.4 & 186 & 3800 & female & 2007\\
\hline
Adelie & Torgersen & 40.3 & 18.0 & 195 & 3250 & female & 2007\\
\hline
Adelie & Torgersen & NA & NA & NA & NA & NA & 2007\\
\hline
Adelie & Torgersen & 36.7 & 19.3 & 193 & 3450 & female & 2007\\
\hline
Adelie & Torgersen & 39.3 & 20.6 & 190 & 3650 & male & 2007\\
\hline
\end{tabular}
\end{table}

Figure \ref{fig:penguins-ggplot} shows an plot of the penguins data, made using the \CRANpkg{ggplot2} package.

\begin{verbatim}
penguins %>% 
  ggplot(aes(x = bill_depth_mm, y = bill_length_mm, 
             color = species)) + 
  geom_point()
\end{verbatim}

\begin{figure}
\includegraphics[width=1\linewidth]{paper-blocking_files/figure-latex/penguins-ggplot-1} \caption{A basic non-interactive plot made with the ggplot2 package on palmer penguin data. Three species of penguins are plotted with bill depth on the x-axis and bill length on the y-axis. Visit the online article to access the interactive version made with the plotly package.}\label{fig:penguins-ggplot}
\end{figure}

\section{Summary}\label{summary}

We have displayed various tooltips that are available in the package \pkg{ToOoOlTiPs}.

\section{Acknowledgements}\label{acknowledgements}

Work on this package is supported by the National Science Centre, OPUS 20 grant no. 2020/39/B/HS4/00941

\bibliography{RJreferences.bib}

\address{%
Maciej Beręsewicz\\
University of Economics and BusinessStatisical Office in Poznań\\%
Department of Statistics, Poznań, Poland\\ Centre for the Methodology of Population Studies\\
%
\url{https://maciejberesewicz.com}\\%
\textit{ORCiD: \href{https://orcid.org/0000-0002-8281-4301}{0000-0002-8281-4301}}\\%
\href{mailto:maciej.beresewicz@poznan.pl}{\nolinkurl{maciej.beresewicz@poznan.pl}}%
}

\address{%
Adam Struzik\\
Adam Mickiewicz UniversityStatisical Office in Poznań\\%
Department of Mathematics, Poznań, Poland\\ Centre for Urban Statistics\\
%
%
%
\href{mailto:adastr5@st.amu.edu.pl}{\nolinkurl{adastr5@st.amu.edu.pl}}%
}
